% ============================================================
% CAPITOLUL 3 — TEHNOLOGII ȘI INSTRUMENTE UTILIZATE
% ============================================================

\chapter{Tehnologii și Instrumente Utilizate}
\label{ch:tehnologii}

\section{Limbaje de programare pentru calcul științific}
\label{sec:limbaje}

Alegerea limbajului de programare pentru implementarea modelelor matematice
nu este neutră. Performanța, expresivitatea și ecosistemul de biblioteci
influențează atât timpul de dezvoltare, cât și comportamentul numeric al aplicației.
În această lucrare, R, Julia și Python sunt tratate ca ecosisteme de calcul
științific, nu doar ca sintaxe diferite pentru aceleași formule.

% ---- R ---------------------------------------------------
\subsection{R}
\label{subsec:r}

\subsubsection{Istoric și filosofie}

R este un limbaj de programare și mediu software pentru calcul statistic și
grafică, descendent direct al limbajului S (Bell Laboratories, 1976). A fost
creat de Ross Ihaka și Robert Gentleman la Universitatea din Auckland în 1993,
iar în 1997 a fost adoptat ca proiect open-source de \emph{R Development Core
Team} \cite{rcoreteam2024}. Filosofia R este centrată pe date: vectorii,
matricile și cadrele de date (data frames) sunt cetățeni de primă clasă,
iar gramatica limbajului este proiectată pentru manipularea și analiza datelor.

\subsubsection{Caracteristici tehnice relevante}

R este un limbaj \emph{interpretat}, dinamic tipizat, cu suport pentru
programare funcțională, procedurală și orientată pe obiecte prin sistemele S3,
S4 și R6. Pentru calcul numeric, avantajul său principal vine din vectorizarea
nativă: operațiile pe vectori și matrice sunt executate frecvent în cod C sau
Fortran, evitând buclele interpretate. În același timp, modelul de memorie cu
garbage collector și copiere la scriere poate introduce costuri vizibile atunci
când sunt manipulate structuri mari. Posibilitatea de a apela cod nativ prin
\texttt{.C()}, \texttt{.Call()} și \texttt{.Fortran()} explică de ce multe pachete
R performante ascund implementări compilate în spatele unei interfețe statistice
convenabile.

\subsubsection{Pachete relevante pentru modele compartimentale}

Pentru modelele compartimentale implementate, pachetul relevant este \texttt{deSolve}
\cite{soetaert2010desolve}, care oferă funcția \texttt{ode()} și acces la metode
consacrate precum \texttt{lsoda}, \texttt{ode45} și \texttt{rk4}. În jurul lui se
pot folosi pachete precum \texttt{ggplot2} pentru vizualizare, \texttt{data.table}
pentru prelucrare tabulară și \texttt{bench} sau \texttt{microbenchmark} pentru
măsurarea timpilor de execuție. În implementarea sistemului, runner-ul R folosește
\texttt{deSolve} pentru integrarea numerică, iar vizualizarea este lăsată aplicației
web.

% ---- JULIA -----------------------------------------------
\subsection{Julia}
\label{subsec:julia}

\subsubsection{Istoric și filosofie}

Julia a fost creat la MIT de Jeff Bezanson, Alan Edelman, Stefan Karpinski
și Viral B. Shah, cu versiunea 1.0 lansată în august 2018 \cite{bezanson2017julia}.
Motivația proiectului a fost depășirea \emph{problemei celor două limbaje}
(two-language problem): cercetătorii scriu prototipuri în R/Python (productivitate
ridicată, viteză scăzută) și re-implementează codul critic în C/Fortran (viteză
ridicată, productivitate scăzută). Julia vizează ambele obiective simultan.

\subsubsection{Compilare Just-In-Time și sistemul de tipuri}

Julia utilizează compilare \textbf{JIT} (Just-In-Time) bazată pe LLVM, cu
o particularitate fundamentală: specializarea pe tipuri concrete. Când o funcție
este apelată pentru prima dată cu argumente de tipuri specifice, compilatorul
generează cod mașină optimizat pentru acele tipuri exacte. Apelurile ulterioare
cu aceleași tipuri rulează direct codul compilat, apropiind performanța de C.

Mecanismul de \emph{multiple dispatch} permite definirea mai multor metode pentru
aceeași funcție, în funcție de tipurile argumentelor. Împreună cu sistemul de tipuri
parametric, acesta permite compilatorului să genereze cod specializat pentru cazurile
folosite efectiv. Macrocomenzile completează acest model prin posibilitatea de a
transforma codul la momentul parsării, o tehnică utilizată frecvent în ecosistemul
științific Julia.

\subsubsection{DifferentialEquations.jl}

Ecosistemul \textbf{SciML} (Scientific Machine Learning) al Julia conține
pachetul \texttt{DifferentialEquations.jl} \cite{rackauckas2017diffeq}, un
framework open-source matur pentru rezolvarea ecuațiilor diferențiale.
Pachetul include solver-e pentru ecuații diferențiale ordinare, parțiale,
stochastice și cu întârziere. Pentru probleme nestiff, metode precum
\texttt{Tsit5} și \texttt{Vern7/9} sunt folosite frecvent, iar pentru probleme
stiff există familii precum \texttt{Rodas}. Ecosistemul SciML integrează și
diferențiere automată, ceea ce îl face relevant pentru estimarea parametrilor și
identificarea sistemelor dinamice.

% ---- PYTHON ----------------------------------------------
\subsection{Python}
\label{subsec:python}

\subsubsection{Istoric și filosofie}

Python a fost creat de Guido van Rossum și lansat în 1991 \cite{vanrossum2009python}.
Filosofia sa, rezumată în \emph{Zen of Python} (\texttt{PEP 20}), accentuează
lizibilitatea și simplitatea codului. Python a devenit limbajul dominant în
știința datelor și machine learning datorită ecosistemului său extins și
comunității masive.

\subsubsection{Caracteristici tehnice și ecosistem științific}

Python este un limbaj interpretat, dinamic tipizat, multi-paradigmă. Principala
limitare pentru calcul intensiv este \textbf{GIL} (Global Interpreter Lock),
care previne execuția paralelă a thread-urilor Python pure. Performanța este
recuperată prin biblioteci care rulează cod nativ C sau Fortran în afara GIL.
NumPy oferă calcul vectorizat pe array-uri dense, SciPy
\cite{virtanen2020scipy} adaugă algoritmică științifică și modulul
\texttt{scipy.integrate}, iar Pandas este folosit pe scară largă pentru date
tabulare. Pentru sistemul dezvoltat aici, componenta esențială este
\texttt{solve\_ivp}, care expune metode precum RK45, DOP853, Radau, BDF și LSODA.

\subsubsection{Funcțiile de rezolvare a EDO în SciPy}

Interfața \texttt{solve\_ivp} din \texttt{scipy.integrate} este apelată printr-o
funcție RHS care descrie modelul selectat. Exemplul următor folosește cazul SIR,
deoarece are doar trei variabile de stare; pentru SIDARTHE, aceeași
interfață primește un vector de opt compartimente și un dicționar mai
mare de parametri.
\begin{lstlisting}[language=Python, caption={Interfața solve\_ivp din SciPy pentru un sistem compartimental SIR}]
from scipy.integrate import solve_ivp
import numpy as np

def sir_model(t, y, beta, gamma, N):
    S, I, R = y
    dSdt = -beta * S * I / N
    dIdt =  beta * S * I / N - gamma * I
    dRdt =  gamma * I
    return [dSdt, dIdt, dRdt]

# Parametri
N = 1_000_000; beta = 0.3; gamma = 0.1
S0, I0, R0 = N - 1, 1, 0
t_span = (0, 160); t_eval = np.linspace(0, 160, 1601)

sol = solve_ivp(
    fun=lambda t, y: sir_model(t, y, beta, gamma, N),
    t_span=t_span,
    y0=[S0, I0, R0],
    method='RK45',
    t_eval=t_eval,
    rtol=1e-8,
    atol=1e-10
)
\end{lstlisting}

% ---- COMPARAȚIE -----------------------------------------
\subsection{Analiză comparativă a limbajelor}
\label{subsec:comparatie_limbaje}

\Cref{tab:comparatie_limbaje} sintetizează principalele caracteristici ale
celor trei limbaje din perspectiva utilizării pentru implementarea modelelor
compartimentale implementate în sistem.

\begin{table*}[t]
\centering
\caption{Comparație a limbajelor de programare utilizate}
\label{tab:comparatie_limbaje}
\renewcommand{\arraystretch}{1.3}
\begin{tabular}{lccc}
\toprule
\textbf{Criteriu} & \textbf{R} & \textbf{Julia} & \textbf{Python} \\
\midrule
Licență          & Open-source & Open-source & Open-source \\
Paradigmă        & Funcț./OOP  & Multi       & Multi       \\
Compilare        & Interpretat & JIT (LLVM)  & Interpretat \\
Solver EDO       & deSolve     & DiffEq.jl   & SciPy       \\
Performanță rel. & Medie       & Înaltă      & Medie       \\
Maturitate ecos. & Înaltă      & Medie       & Înaltă      \\
Utilizare acad.  & Ridicată    & Medie       & Ridicată    \\
\bottomrule
\end{tabular}
\end{table*}

% ============================================================
\section{Baze de date distribuite — YugabyteDB}
\label{sec:yugabytedb}
% ============================================================

\subsection{Concepte de baze de date distribuite}
\label{subsec:baze_distribuite}

O \emph{bază de date distribuită} este o colecție de date distribuite fizic
pe mai multe noduri (mașini sau locații geografice), gestionată de un sistem
software care oferă transparența localizării: aplicațiile interacționează cu
baza de date ca și cum ar fi centralizată.

\begin{definitie}[Teorema CAP]
\label{def:cap}
Teorema CAP (Brewer, 2000) afirmă că un sistem distribuit nu poate garanta
simultan mai mult de două dintre următoarele trei proprietăți:
\emph{consistency}, adică nodurile observă aceeași stare logică a datelor;
\emph{availability}, adică cererile valide primesc răspuns chiar în condiții
de degradare; și \emph{partition tolerance}, adică sistemul continuă să funcționeze
în prezența unor întreruperi de comunicare între noduri.
\end{definitie}

Bazele de date distribuite precum YugabyteDB, CockroachDB și Spanner optează
pentru \textbf{CP} (consistency + partition tolerance), sacrificând disponibilitatea
perfectă în favoarea coerenței datelor.

\subsection{Arhitectura YugabyteDB}
\label{subsec:yugabyte_arhitectura}

YugabyteDB \cite{yugabyte2024} este o bază de date distribuită open-source,
construită pentru a combina caracteristicile PostgreSQL cu scalabilitatea
orizontală și toleranța la erori a sistemelor NoSQL. Arhitectura sa se bazează
pe doi piloni principali:

\textbf{1. Stratul de stocare DocDB.} DocDB este motorul de stocare
distribuit al YugabyteDB, bazat pe arhitectura \textbf{LSM-tree} (Log-Structured
Merge-tree) derivată din \textbf{RocksDB}. Datele sunt stocate în \emph{tablete}
--- unități de sharding distribuite pe noduri. Fiecare tabletă este replicată
pe cel puțin trei noduri prin protocolul Raft. Scrierile trec printr-un lider,
sunt replicate către nodurile urmăritoare și sunt confirmate după ce o majoritate
de replici a persistat modificarea.

\textbf{2. Nivelurile de interfață API.} Interfața \textbf{YSQL} este compatibilă
cu PostgreSQL și suportă SQL, tranzacții ACID distribuite, join-uri și indecși.
Interfața \textbf{YCQL} este compatibilă cu Apache Cassandra și este orientată
spre workload-uri wide-column. În această lucrare este folosită interfața YSQL,
pentru compatibilitatea cu driver-ele PostgreSQL disponibile în Go, R, Julia și
Python.

\subsection{Sharding și replicare}
\label{subsec:sharding}

YugabyteDB suportă sharding prin hash și sharding pe intervale. Hash sharding
distribuie uniform cheia primară pe tablete și reduce riscul de noduri fierbinți,
dar este mai puțin potrivit pentru scanări ordonate pe intervale. Range sharding
păstrează vecinătatea valorilor și este util pentru serii temporale, însă poate
concentra scrierile dacă toate valorile noi ajung în același interval.

\begin{observatie}
Pentru datele epidemiologice cu interogări frecvente pe intervale de date
(serii temporale), range sharding pe coloana \texttt{timestamp} sau
\texttt{date} este strategia recomandată, împreună cu un index secundar
pe codul bolii și/sau regiune geografică.
\end{observatie}

% ============================================================
\section{Orchestrarea containerelor — Kubernetes}
\label{sec:kubernetes}
% ============================================================

\subsection{Containere și nevoia de orchestrare}
\label{subsec:containere}

Containerizarea (Docker) permite împachetarea unei aplicații împreună cu
toate dependențele sale într-o unitate portabilă și reproductibilă.
Pe măsură ce sistemele cresc în complexitate și numărul de containere
depășește câteva zeci, gestionarea manuală devine infezabilă: unde rulează
fiecare container? Cum se reporneşte dacă cade? Cum se scalează?

\textbf{Kubernetes} (K8s) \cite{burns2016kubernetes, kubernetes2024} este
platforma open-source de orchestrare a containerelor, inițial dezvoltată
la Google pe baza experienței cu sistemul intern Borg, și donată CNCF
(Cloud Native Computing Foundation) în 2014.

\subsection{Arhitectura Kubernetes}
\label{subsec:k8s_arhitectura}

Un cluster Kubernetes este împărțit între planul de control și nodurile worker.
Planul de control expune API-ul clusterului prin \texttt{kube-apiserver}, păstrează
starea în \texttt{etcd}, programează pod-urile prin \texttt{kube-scheduler} și
rulează controllerele care aduc starea curentă spre starea dorită. Pe nodurile
worker, \texttt{kubelet} gestionează ciclul de viață al pod-urilor,
\texttt{kube-proxy} configurează accesul de rețea la servicii, iar container
runtime-ul pornește efectiv containerele.

\subsection{Resurse Kubernetes relevante}
\label{subsec:k8s_resurse}

Resursele folosite cel mai des sunt \texttt{Pod}, \texttt{Deployment},
\texttt{StatefulSet}, \texttt{Service}, \texttt{PersistentVolumeClaim},
\texttt{ConfigMap} și \texttt{Secret}. Pentru API și frontend-uri sunt suficiente
Deployment-uri, deoarece instanțele nu păstrează stare locală. Pentru YugabyteDB
este necesar un StatefulSet, care oferă identitate stabilă și volume persistente
pentru fiecare replică.

% ============================================================
\section{Backend API — Go}
\label{sec:golang}
% ============================================================

\subsection{Caracteristici și filosofia limbajului Go}
\label{subsec:go_caracteristici}

Go (cunoscut și ca Golang) este un limbaj de programare open-source creat
la Google de Robert Griesemer, Rob Pike și Ken Thompson, cu prima versiune
stabilă în 2012 \cite{golangspec2024}. Go a fost proiectat pentru productivitate
la scară mare, cu accent pe simplitate, compilare rapidă și concurență structurată.

Pentru un API backend, Go are câteva proprietăți utile: compilează rapid la cod
mașină, produce binare statice ușor de containerizat, oferă concurență prin
goroutine și canale, iar biblioteca standard include un pachet HTTP suficient de
solid pentru servicii REST. Interfețele structurale reduc cuplarea dintre module,
iar garbage collectorul concurrent menține latențe acceptabile pentru un API care
nu face calcul numeric intensiv în procesul HTTP.

\subsection{Ecosistem pentru API-uri REST}
\label{subsec:go_api}

Deși \texttt{net/http} din biblioteca standard este suficient pentru API-uri simple,
proiectul folosește routerul \texttt{chi}, o bibliotecă idiomatică și redusă ca
suprafață, care oferă rutare ierarhică, middleware-uri compozabile și integrare
directă cu tipurile standard \texttt{http.Handler}. Această alegere evită un
framework mai greu și se potrivește cu dimensiunea backend-ului implementat.

% ============================================================
\section{Frontend — Vue.js}
\label{sec:vuejs}
% ============================================================

\subsection{Arhitectura Vue 3 și Composition API}
\label{subsec:vue_arhitectura}

Vue.js \cite{vuejs2024} este un framework JavaScript progresiv pentru construirea
interfețelor utilizator, creat de Evan You în 2014. Vue 3 (2020) introduce
\textbf{Composition API}, o alternativă la Options API care permite organizarea
logicii prin funcționalitate (nu prin tip --- data, methods, computed) și
reutilizarea logicii prin \emph{composables}.

Arhitectura Vue 3 combină un sistem reactiv bazat pe \texttt{Proxy}, un Virtual DOM
pentru reconcilierea arborelui de interfață și Single File Components, în care
template-ul, logica TypeScript și stilurile CSS sunt grupate în fișiere
\texttt{.vue}. Pentru aplicația de față, această structură permite păstrarea
componentelor de formular, navigare și grafic în unități ușor de urmărit.

\subsection{Ecosistem și biblioteci utilizate}
\label{subsec:vue_ecosistem}

În implementare sunt folosite Vue Router pentru rutare client-side, Pinia pentru
starea globală, Axios pentru apelurile REST și Chart.js, prin \texttt{vue-chartjs},
pentru reprezentarea curbelor compartimentale ale modelului selectat. Componenta
de grafic construiește dinamic seriile în funcție de cheile prezente în obiectul
JSONB \texttt{compartments}, astfel încât același cod afișează atât traiectoriile
$S(t)$, $I(t)$, $R(t)$ ale modelului SIR, cât și cele opt traiectorii ale
modelului SIDARTHE. Aceste biblioteci acoperă nevoile aplicației fără a introduce
un framework frontend suplimentar peste Vue.
